\capitulo{1}{Introducción}


Las enfermedades neuromotoras engloban un conjunto de trastornos que afectan al control del movimiento, provocando en muchos casos la aparición de movimientos involuntarios. En pacientes que padecen este tipo de patologías, como la parálisis cerebral atetoide, una de las principales dificultades consiste en distinguir entre movimientos voluntarios e involuntarios. Esta diferenciación resulta fundamental tanto para la evaluación del paciente como para la interpretación de sus respuestas al interaccionar con diferentes profesionales. 

En este contexto, surge la necesidad de desarrollar herramientas que permitan analizar de la manera más objetiva posible los movimientos. El presente trabajo propone el diseño de soluciones tecnológicas orientadas a la detección y análisis de movimientos involuntarios, centrando su aplicación en pacientes con parálisis cerebral atetoide. Para ello, se abordan los conceptos teóricos de la patología y se plantean posibles soluciones que puedan facilitar la evaluación y detección de estos movimientos. 






